% 67-50 ①(67쪽) — 보기 그래프: x=a 에서 접선이 수평, x=b 에서 극소, x=c 에서 극대인 개형.
% 그림 자체가 보기라 TikZ 로 다시 그림. 라벨은 원문 그대로 a · b · c · O · x · y 만.
\begin{tikzpicture}[x=0.46cm, y=0.46cm, line width=0.5pt, baseline=(current bounding box.center)]
  \path (-2.35,-1.15) rectangle (2.95,3.00);
  \draw[->, >=stealth, line width=0.7pt] (-2.15,0) -- (2.80,0) node[below=1pt, font=\small] {$x$};
  \draw[->, >=stealth, line width=0.7pt] (0,-0.95) -- (0,2.85) node[left=1pt, font=\small] {$y$};
  \draw[densely dashed, line width=0.4pt] (-1.0,0) -- (-1.0,1.28);
  \draw[densely dashed, line width=0.4pt] (0.8,0) -- (0.8,0.38);
  \draw[densely dashed, line width=0.4pt] (1.8,0) -- (1.8,1.65);
  \draw[line width=0.6pt] plot[smooth, tension=0.7] coordinates {
    (-1.90,2.50) (-1.65,1.95) (-1.45,1.60) (-1.25,1.38) (-1.00,1.28) (-0.75,1.18)
    (-0.50,1.00) (-0.20,0.75) (0.10,0.55) (0.40,0.42) (0.80,0.38) (1.10,0.55)
    (1.35,0.95) (1.60,1.45) (1.80,1.65) (2.00,1.50) (2.20,1.00) (2.40,0.30) (2.55,-0.30)};
  \node[below=1pt, font=\small] at (-1.0,0) {$a$};
  \node[above left=-2pt and -3pt, font=\small] at (0,0) {$\pt{O}$};
  \node[below=1pt, font=\small] at (0.8,0) {$b$};
  \node[below=1pt, font=\small] at (1.8,0) {$c$};
\end{tikzpicture}
